% Generated draft: requires mathematical and source review.
\documentclass[11pt]{article}
\usepackage[margin=25mm]{geometry}
\usepackage{fontspec}
\setmainfont{Latin Modern Roman}
\usepackage{amsmath,amssymb,mathtools,unicode-math}
\setmathfont{Latin Modern Math}
\usepackage{enumitem,needspace,xurl,hyperref}
\setlength{\parindent}{0pt}
\setlength{\parskip}{5pt}
\setlength{\emergencystretch}{3em}
\newcommand{\N}{\mathbb{N}}
\newcommand{\Z}{\mathbb{Z}}
\newcommand{\Q}{\mathbb{Q}}
\newcommand{\R}{\mathbb{R}}
\newcommand{\C}{\mathbb{C}}
\newcommand{\F}{\mathbb{F}}
\newcommand{\A}{\mathbb{A}}
\newcommand{\PP}{\mathbb{P}}
\newcommand{\E}{\mathbb{E}}
\newcommand{\eps}{\varepsilon}
\newcommand{\dd}{\,\mathrm{d}}
\begin{document}
\section*{UnsolvedMath ID: 30002219}
\textbf{Source problem code:} catalogue entry
\subsection{Definitions and mathematical statement}
Let \(A\) and \(B\) be arbitrary separable complex \(C^*\)-algebras; neither is assumed unital. Their direct sum \(A\oplus B\) has coordinatewise algebraic operations and involution, with \(\|(a,b)\|=\max\{\|a\|,\|b\|\}\).

The supplied record does not define generator rank. To make the question explicit, use the following unitization convention, which requires confirmation against the cited source. For a unital \(C^*\)-algebra \(D\), write \(D_{\mathrm{sa}}=\{d\in D:d=d^*\}\), and define
\[
\operatorname{Gen}_{m}(D)_{\mathrm{sa}}=\{(d_1,\ldots,d_m)\in D_{\mathrm{sa}}^m:C^*(d_1,\ldots,d_m)=D\},\qquad m\geq 1.
\]
Here \(C^*(d_1,\ldots,d_m)\) is the smallest norm-closed \(*\)-subalgebra containing the listed elements; the identity is not adjoined automatically. Set
\[
\operatorname{gr}(D)=\min\bigl\{k\in\{0,1,2,\ldots\}:\operatorname{Gen}_{k+1}(D)_{\mathrm{sa}}\text{ is dense in }D_{\mathrm{sa}}^{k+1}\bigr\},
\]
with value \(\infty\) if this set is empty, and with density taken in the maximum product norm. For nonunital \(D\), set \(\operatorname{gr}(D)=\operatorname{gr}(\widetilde D)\), where \(\widetilde D\) is its minimal unitization.

Under this explicit convention, the question is whether, for every such pair \(A,B\),
\[
\operatorname{gr}(A\oplus B)=\max\{\operatorname{gr}(A),\operatorname{gr}(B)\},
\]
where the maximum is taken in \(\{0,1,2,\ldots\}\cup\{\infty\}\).

The supplied dated curation labels the problem partially solved and reports the formula as proved for commutative, real-rank-zero, and subhomogeneous summands, while describing the arbitrary separable case as unclear. These are supplied literature qualifications, not independently verified conclusions; the record does not specify the precise hypotheses of those reported special cases.
\subsection{Short English statement}
For every pair of separable complex \(C^*\)-algebras, does taking their direct sum make the generator rank equal to the larger of the two individual ranks? The draft explicitly uses the unitization convention for nonunital algebras, pending source confirmation. The supplied partially solved label and special-case claims are dated curation requiring independent review.
\subsection{Sources}
\begin{itemize}
\item Supplied local UnsolvedMath registry record: permanent ID \(30002219\), title \emph{Generator Rank of Direct Sums}, source problem code \texttt{catalogue entry}, in the Oberwolfach Reports Open Problems collection. No individual proposer is named; the supplied proposed year is 2012.
\item Supplied source citation: \emph{C*-Algebras, Dynamics, and Classification} (2013), Oberwolfach Reports, pp.\ 3129--3209, DOI 10.4171/owr/2012/52. \url{https://doi.org/10.4171/owr/2012/52}. This is also the supplied literature link labelled ``original OWR question.''
\item Supplied literature link labelled ``generator rank of C-star algebras'': \url{https://doi.org/10.1016/j.jfa.2020.108874}.
\item Supplied literature link labelled ``subhomogeneous direct sums'': \url{https://www.cambridge.org/core/journals/canadian-journal-of-mathematics/article/generator-rank-of-subhomogeneous-calgebras/976DD3CC21BFD2F7688383FC08C8EE9C}.
\item Literature assessment dated exactly \texttt{2026-08-21}: the supplied curation reports the direct-sum formula for commutative, real-rank-zero, and subhomogeneous summands and describes the general separable case as unclear. This is supplied dated curation derived from the OWR record, not an independently verified literature assessment.
\end{itemize}
\end{document}
